\newpage
\section{Komentarze w postaci obrazków (3.5)}
\label{sec:obrazy}

Na forach część treści to memy i zrzuty ekranu, których zwykły klasyfikator tekstu nie przeczyta. Dlatego ocena
obrazka ma dwa kroki. Najpierw model wizyjny \kod{gemma3:4b} dostaje obrazek i polecenie: ,,przepisz cały tekst
z obrazka i jednym zdaniem opisz, co przedstawia''. Odpowiedź ma ustalony układ (\kod{TEXT:} i \kod{PICTURE:}), z którego
wycinam tekst. Potem ten tekst przechodzi przez tę samą ocenę co zwykły komentarz (domyślnie w trybie RAG).

Mały model wizyjny czasem powtarzał tę samą linię dwa razy albo dopisywał etykietę \kod{TEXT:} w środku odpowiedzi.
Funkcja \kod{clean\_text} usuwa etykiety i powtórzone linie, a następnie skleja tekst w jedno zdanie.

\zrzut{code_vision}{Kod: odczyt tekstu z obrazka i ocena (\kod{vision.py}).}

Obrazki testowe wygenerowałem skryptem \kod{scripts/make\_sample\_images.py} (Pillow): trzy memy z tekstem na górze
i na dole oraz trzy zrzuty ekranu wpisu z forum. Dzięki temu nie używam cudzych materiałów.

\begin{figure}[H]
  \centering
  \includegraphics[width=0.45\linewidth]{meme_03.png}\hfill
  \includegraphics[width=0.52\linewidth]{screenshot_01.png}
  \caption{Przykładowe obrazki: mem z groźbą i zrzut ekranu ze spamem.}
\end{figure}

\zrzut{image_meme}{Wynik: \kod{moderator image data/images/meme\_03.png}.}

\zrzut{image_screenshot}{Wynik: \kod{moderator image data/images/screenshot\_01.png}.}

Model wizyjny przeczytał oba obrazki poprawnie, z jednym błędem: na memie ,,NAPISZ'' odczytał jako ,,NAPIS''.
Dla oceny nie miało to znaczenia -- komentarz dostał kategorię \kod{threat}, stopień 3 i paragraf §3.1, a reklama ze
zrzutu ekranu kategorię \kod{spam} i paragraf §5.1. Ocena obrazka trwa około 16--18 sekund, bo do pamięci trzeba
załadować dwa modele.
